% cv-body.tex — content shared by the short CV and the full dossier.
% \ifdossier ... \else ... \fi selects dossier-only or CV-only material.

\cvheader{\cvname, Ph.D.}
  {Professor, Department of Electrical and Electronic Engineering, BUET, Dhaka}
  {\cvcontact{\faEnvelope[regular]}{mailto:sajid@eee.buet.ac.bd}{sajid@eee.buet.ac.bd}\cvsep
   \cvcontact{\faGlobe}{https://www.sajid.bd}{sajid.bd}\cvsep
   \cvcontact{\aiOrcid}{https://orcid.org/0000-0002-0216-7125}{0000-0002-0216-7125}\cvsep
   \cvcontact{\aiGoogleScholar}{https://scholar.google.com/citations?user=Fu8Hkb4AAAAJ}{Google Scholar}\cvsep
   \cvcontact{\faLinkedin}{https://www.linkedin.com/in/sajidmc/}{sajidmc}\cvsep
   \cvcontact{\faGithub}{https://github.com/sajidmc}{sajidmc}}

% ═══════════════════════════════════════════════════════════════════════════
\section{Profile}
\begin{cvblock}
Professor in the Department of Electrical and Electronic Engineering at the Bangladesh University of Engineering and Technology (BUET); Ph.D., Purdue University. Research spans quantum computing, photonics, antennas, computing, embedded systems and renewable energy. Senior Member of IEEE, Founding President of the Optica Bangladesh Section, and Founding Chair of the IEEE Photonics Society Bangladesh Chapter.
\end{cvblock}

% ═══════════════════════════════════════════════════════════════════════════
\section{Appointments}
\cvgroup{Bangladesh University of Engineering and Technology (BUET), Dhaka}
\cvitem{Jul 2025 -- present}{\textbf{Professor}, Department of Electrical and Electronic Engineering}
\cvitem{Jul 2022 -- Jul 2025}{\textbf{Associate Professor}, Department of Electrical and Electronic Engineering}
\cvitem{Jun 2013 -- Jul 2022}{\textbf{Assistant Professor}, Department of Electrical and Electronic Engineering}
\cvitem{Jan 2010 -- Jun 2013}{\textbf{Lecturer}, Department of Electrical and Electronic Engineering}
\cvitem{Nov 2009 -- Jan 2010}{\textbf{Lecturer}, Institute of Information and Communication Technology (IICT)}

% ═══════════════════════════════════════════════════════════════════════════
\section{Education}
\cvitem{Aug 2013 -- Aug 2019}{\textbf{Ph.D.}, School of Electrical and Computer Engineering, Purdue University, West Lafayette, IN, USA\\
  \cvmeta{Thesis: \emph{Wavefront Manipulation with Metasurfaces Based on New Materials}.
  Co-supervisors: Alexandra Boltasseva and Alexander V. Kildishev.}}
\cvitem{Aug 2011 -- 2013}{\textbf{M.Sc.\ Engg.}, Department of Electrical and Electronic Engineering, BUET\\
  \cvmeta{Thesis: \emph{Design of a Fractal Antenna Based on Hexaflake Fractal Structure}. Supervisor: Dr.\ M.\,A. Matin.}}
\cvitem{Dec 2004 -- Aug 2010}{\textbf{B.Sc.\ Engg.}, Department of Electrical and Electronic Engineering, BUET\\
  \cvmeta{CGPA 3.94/4.00. Thesis: \emph{Design and Analysis of a Multiband Dual Feed Axially Symmetric Cassegrain Antenna System}. Supervisor: Dr.\ M.\,A. Matin.}}
\ifdossier
\cvitem{2004}{\textbf{H.S.C.}, Notre Dame College, Dhaka \cvmeta{\enspace GPA 5.00/5.00}}
\cvitem{2002}{\textbf{S.S.C.}, Udayan Uchchya Madhyamic Bidyalaya, Dhaka \cvmeta{\enspace GPA 5.00/5.00}}
\fi

% ═══════════════════════════════════════════════════════════════════════════
\section{Research}
\begin{cvblock}
Research is organised under the \textbf{Q-PACERS} framework: Quantum, Photonic, Antenna, Computing, Embedded and Renewable-energy Systems, within the Electronics and Photonics (EP) division of EEE.
\ifdossier\else
\begin{cvlist}
  \item \textbf{Nanophotonics and plasmonics:} plasmonic nanostructures for optical sensing and modulation; dual-band near- and mid-infrared absorber for biochemical sensing; plasmonic biosensors and electro-optic modulators.
  \item \textbf{Metasurfaces and antennas:} fractal (hexaflake) microstrip antennas; electrically reconfigurable metasurfaces for near-infrared metalensing; metasurfaces for imaging and holography.
  \item \textbf{Quantum computing and quantum photonics:} quantum algorithms, photonic qubits and integrated photonic circuits, and loss-tolerant one-way quantum repeaters.
  \item \textbf{Embedded systems and IoT:} sensor-driven embedded platforms for environmental monitoring, smart grids and healthcare.
  \item \textbf{Renewable energy:} light trapping in heterojunction solar cells with hybrid metal--dielectric nanostructures; photocatalytic water splitting.
\end{cvlist}
\fi
\end{cvblock}
\ifdossier
\begin{cvblock}
\textbf{Nanophotonics and plasmonics.} The group develops photonic materials and devices at the nanoscale, designing plasmonic nanostructures for optical sensing and modulation. A recent publication reports a concentric plasmonic resonator for dual-band absorption in the near- and mid-infrared, aimed at biochemical sensing. Plasmonic biosensors and high-speed electro-optic modulators are continuing interests, with the goal of nanophotonic systems that extend light--matter interaction for communications, healthcare and computing.

\textbf{Advanced antennas and metasurfaces.} Building on a background in antenna engineering (his M.Sc.\ research produced a microstrip patch antenna with a hexaflake fractal geometry), the group studies metasurfaces, engineered electromagnetic surfaces, for new applications. A recent thesis student developed an electrically reconfigurable metasurface for metalensing in the near-infrared band. Metasurface-based antennas and optical elements are central to this programme, with applications from wireless communication devices to compact imaging and holography.

\textbf{Quantum computing and quantum photonics.} The research extends to quantum information technology: developing and optimising quantum algorithms and exploring practical quantum hardware, in particular photonic systems (photons in nanoscale devices and integrated photonic circuits) that realise qubits and quantum logic. A dedicated graduate course on quantum computing and photonic technologies (see Teaching) lays the groundwork for a research programme in quantum algorithms and photonic quantum hardware, contributing to Bangladesh's entry into quantum technology research.

\textbf{Embedded systems and computing for IoT.} Design projects pursue efficient computing systems and embedded platforms for the Internet of Things. By integrating sensors, processors and communication interfaces, these platforms are tailored to environmental monitoring, smart grids and healthcare devices. The work complements EEE 416 (Microprocessor and Embedded Systems Design Laboratory), where undergraduate teams prototype IoT systems that contribute to edge-computing and distributed-sensor-network research.

\textbf{Renewable-energy and sustainable systems.} The group works on efficient solar-cell design and grid optimisation, and on improving photovoltaic devices through nanotechnology. A recent journal paper demonstrates an efficiency improvement in a crystalline-silicon/thin-film solar cell using hybrid metal--dielectric nanostructures. Photocatalytic water splitting for hydrogen production is a further interest. The plan is to integrate these strands into smart renewable-energy systems, for instance solar converters with embedded IoT sensors for performance optimisation, or photonic devices that improve energy generation and storage.
\end{cvblock}
\fi

% ═══════════════════════════════════════════════════════════════════════════
\section{Research Grants}
\cvitem{Apr 2024 -- Oct 2025}{\textbf{Synergizing Deep Learning and Topology Optimization for Tunable Meta-lens Design}\\
  \cvmeta{Principal Investigator. Internal Research Grant, BUET. BDT 14,00,000.}}
\cvitem{Jan 2023 -- Jan 2024}{\textbf{Structurally Tunable Gear-Shaped Plasmonic Sensor}\\
  \cvmeta{Principal Investigator. Basic Research Grant, Research and Innovation Centre for Science and Engineering (RISE), BUET. BDT 7,00,000.}}
\cvitem{Jan 2021 -- Jan 2022}{\textbf{Design and Low-cost Fabrication of Optical Biosensors for Virus Detection}\\
  \cvmeta{Principal Investigator. Basic Research Grant, Center for Advanced Study and Research (CASR), BUET. BDT 1,00,000.}}

% ═══════════════════════════════════════════════════════════════════════════
\section{Teaching and Curriculum Development}
\ifdossier
\begin{cvblock}
Course content is updated continually to reflect the state of the art, with an emphasis on hands-on learning. Beyond delivering core EEE courses, the teaching programme introduces new courses and pedagogical approaches that prepare students for rapidly evolving fields.
\end{cvblock}
\fi
\cvitem{EEE 415/416}{\textbf{Modernisation of the Microprocessor and Embedded Systems curriculum.}
  \ifdossier
  EEE 415 (Microprocessors and Embedded Systems) traditionally taught the legacy Intel 8086 microarchitecture and its assembly language, at BUET and at most public and private universities in Bangladesh, whereas comparable courses abroad use modern architectures (MIPS, ARM, RISC-V). He led an overhaul that replaced the 8086-centred syllabus with one built on the ARM architecture and microcontroller design. The revised course covers modern processors and microcontrollers, the ARM Cortex-A and Cortex-M architectures, and high-level design with Verilog. Students learn assembly programming and system design on ARM/RISC-V platforms (stack operations, interrupts, serial interfaces), aligning the curriculum with global standards and the needs of Bangladesh's technology sector.
  \else
  Replaced the legacy Intel 8086 syllabus with ARM Cortex-M microcontroller design, ARM/RISC-V assembly and Verilog-based design.
  \fi}
\cvitem{Lab, 2022}{\textbf{Modernisation of the Microprocessor and Embedded Systems Laboratory} (lab-in-charge).
  \ifdossier
  To support the revised syllabus, he submitted a proposal to the Head of the Department, which was forwarded to and approved by the University Syndicate in January 2022.
  \begin{cvlist}
    \item \textbf{Pedagogical layout.} Three zones (Class Space, Project Space and Storage Space) for unobstructed instruction, team prototyping benches and organised component inventory. Each workstation, shared by two students, integrates an FPGA board, an MCU kit and a development PC.
    \item \textbf{Core hardware.} STM32 Nucleo ARM Cortex-M4 evaluation kits and Digilent Nexys A7-100T (Artix-7) FPGA trainers anchor the digital-design workflow, supported by 150\,MHz digital storage oscilloscopes, programmable power supplies, logic analysers, and sensors and actuators (stepper and servo motors, LCDs, temperature sensors).
    \item \textbf{Rapid prototyping.} A 3D printer for enclosures and a soldering station for PCB prototyping.
    \item \textbf{Budget.} BDT 62.42 lakh in total, with components chosen for low cost and local availability: equipment 26\%, test and measurement 36\%, IT and furniture 24\%, civil and electrical works 14\%; 17 complete workstations and specialised project equipment.
  \end{cvlist}
  \else
  BDT 62.42 lakh; 17 workstations with STM32 Nucleo (ARM Cortex-M4) kits and Digilent Nexys A7 FPGA trainers; approved by the University Syndicate, Jan 2022.
  \fi}
\cvitem{EEE 416}{\textbf{Project-based laboratory.}
  \ifdossier
  He redesigned the EEE 416 laboratory around experiential learning in two phases: guided experiments that reinforce fundamentals, followed by open-ended design projects on EEE 415 topics. Student teams build, for example, an embedded IoT device or a small computer system, applying processor architecture, interfacing and programming, and developing problem-solving and hardware-debugging skills in line with Outcome-Based Education principles. The department later adopted the model for most compulsory sessional courses, replacing the capstone project with ``Culmination of Program Outcomes with Sessional Course Projects''.
  \else
  Guided experiments followed by open-ended design projects; the model was later adopted by the department for most compulsory sessional courses.
  \fi}
\cvitem{EEE 6408, 6505}{\textbf{New postgraduate courses:} EEE 6408 Nano Systems and EEE 6505 Nanophotonics and Plasmonics, both offered for the first time in the department.}
\cvitem{EEE 6002, 2024}{\textbf{Quantum Computing and Quantum Photonics} (special topics).
  \ifdossier
  First offered in 2024, the course gives EEE graduate students a foundation in quantum information science, from qubits and quantum logic gates to quantum algorithms (Deutsch--Jozsa, Grover, Shor) and error correction, with an overview of current quantum hardware platforms. It also integrates photonics, covering photonic quantum computers and continuous-variable quantum computing, filling a gap in the curriculum and preparing students and faculty for research in quantum technologies. Further offerings are planned, potentially including undergraduate electives in quantum technology, nanophotonics or advanced embedded systems.
  \else
  Qubits and gates, quantum algorithms, error correction, and photonic and continuous-variable quantum computing.
  \fi}
\cvitem{EEE 6516}{\textbf{Quantum Computing} (3.0 credit hours, EP division), designed from the special-topics course and approved by the Academic Council.
  \ifdossier\\
  \cvmeta{Syllabus: introduction to quantum computing; postulates of quantum mechanics; qubits and single-qubit gates; operators and quantum measurement; quantum information processing, superposition, entanglement, the no-cloning theorem and quantum teleportation; quantum circuits, controlled operations and measurement, universal quantum gates; quantum algorithms: Deutsch--Jozsa, Grover's and Shor's algorithms, the quantum Fourier transform and phase estimation; density matrices and the Bloch-sphere representation; quantum error correction and fault-tolerant architecture; physical realisations: harmonic oscillator, optical quantum computers, trapped ions, superconducting qubits and new implementation schemes.}
  \fi}

% ═══════════════════════════════════════════════════════════════════════════
\section{Postgraduate Supervision}
\subsection{Completed}
\cvitem{Sep 2025}{Puja Das, M.Sc.\ Engg.\\ \cvmeta{\emph{Efficiency Enhancement of HTL-Free Perovskite/Cadmium-Sulfide Heterojunction Solar Cell Incorporating Light Trapping Structure}}}
\cvitem{Jul 2025}{Ayon Sarker, M.Sc.\ Engg.\\ \cvmeta{\emph{Design of Dual-band Plasmonic Absorber for Biomedical Sensing and Environmental Monitoring}}}
\cvitem{Jul 2025}{Soikot Sarkar, M.Sc.\ Engg.\\ \cvmeta{\emph{Hybrid Metal-Dielectric Nanostructures Integrated Heterojunction Thin Film Solar Cell for Efficiency}}}
\cvitem{May 2025}{Md.\ Asif Hossain Bhuiyan, M.Sc.\ Engg.\\ \cvmeta{\emph{Polarization Insensitive Electrically Reconfigurable Metasurface for Metalensing at Near Infrared Waveband}}}
\cvitem{Jan 2025}{Md.\ Mahfuzul Haque, M.Sc.\ Engg.\\ \cvmeta{\emph{Design of Silicon-Carbide Based Single-Quantum-Well White LED}}}
\cvitem{Jun 2024}{Md.\ Ehsanul Karim, M.Sc.\ Engg.\\ \cvmeta{\emph{Phase Change Material Based Broadband Multifunctional Metasurface for the Visible Range}}}
\cvitem{May 2023}{Shamima Akter Mitu, M.Sc.\ Engg.\\ \cvmeta{\emph{Design of an All-Optical Plasmonic Modulator for Two Micrometer Waveband}}}
\subsection{Ongoing}
\cvitem{since Oct 2024}{Mir Md.\ Aminuzzaman, Ph.D.}
\cvitem{since Oct 2024}{Tanvir Ahmed, M.Sc.\ Engg.}
\cvitem{since Apr 2024}{Zaheer Al Jamee; Md.\ Al Shahriar Shakil; Naved Sadat Yamin; Jayed Hasan Rakib; all M.Sc.\ Engg.}
\cvitem{since Apr 2023}{Abu Md.\ Raihan, M.Sc.\ Engg.}
\cvitem{since Apr 2021}{Purbayan Das, M.Sc.\ Engg.}
\ifdossier
\cvitem{B.Sc.\ thesis}{Addrita Basak Nodi; Ramima Mumtarin; Gopal Gosh (B.Sc.\ Engg., 2021 intake).}
\fi

% ═══════════════════════════════════════════════════════════════════════════
\section{Professional Service and Leadership}
\cvgroup{IEEE}
\cvitem{Apr 2022 -- present}{Vice Chair, IEEE Photonics Society Bangladesh Chapter}
\cvitem{Mar 2021 -- Apr 2022}{Founding Chair, IEEE Photonics Society Bangladesh Chapter}
\cvitem{Jul 2021 -- May 2022}{Secretary, IEEE Bangladesh Section}
\cvitem{Mar 2020 -- Apr 2022}{Chair, IEEE Young Professionals Bangladesh}
\cvitem{Jan 2013 -- Dec 2013}{Chair, IEEE Graduates of the Last Decade (GOLD)}
\ifdossier
\cvitem{Jan 2011 -- Dec 2012}{Vice-Chair, IEEE Graduates of the Last Decade (GOLD)}
\cvitem{Jan 2011 -- Dec 2012}{Student Activities Coordinator, IEEE Bangladesh Section}
\cvitem{Jan 2008 -- Aug 2009}{Chair, IEEE BUET Student Branch}
\cvitem{Jan 2007 -- Dec 2008}{Treasurer, IEEE BUET Student Branch}
\fi
\cvgroup{Optica}
\cvitem{May 2022 -- present}{Founding President, Optica Bangladesh Section}
\cvitem{Jul 2022 -- present}{Founding Moderator, BUET Optical and Photonics Society}
\ifdossier
\cvitem{Jun 2016 -- May 2017}{Treasurer, OSA Purdue Chapter, USA}
\cvgroup{Purdue University}
\cvitem{Jun 2017 -- May 2018}{President, Nanotechnology Student Advisory Council (NSAC)}
\cvitem{Jun 2016 -- May 2017}{Vice-President, Nanotechnology Student Advisory Council (NSAC)}
\cvitem{Jul 2017 -- Jun 2018}{President, Bangladesh Students Association (Purdue BDSA)}
\cvitem{Jul 2015 -- Jun 2017}{Treasurer, Bangladesh Students Association (Purdue BDSA)}
\cvitem{Jun 2015 -- May 2016}{Treasurer, SPIE Purdue Chapter}
\else
\cvitem{2015 -- 2018}{Purdue University: President, Nanotechnology Student Advisory Council (2017--18); President, Bangladesh Students Association (2017--18); officer roles in the OSA and SPIE student chapters.}
\fi

% ═══════════════════════════════════════════════════════════════════════════
\section{Memberships}
\cvitem{Apr 2021 -- present}{Member, National Young Academy of Bangladesh (NYAB)}
\cvitem{}{Senior Member, Institute of Electrical and Electronics Engineers (IEEE); Member, IEEE Photonics Society}
\cvitem{}{Member, Optica}
\cvitem{Apr 2024 -- present}{Life Member, American Alumni Association of Bangladesh (AAAB)}
\cvitem{Apr 2021 -- present}{Life Member, Association of BUET Alumni}

% ═══════════════════════════════════════════════════════════════════════════
\clearpage
\section{Publications}
\cvmetricspanel
\vspace{6pt}
{\small\color{cvgrey}Own name in \textbf{bold}; \textsuperscript{\dag}\,student supervised. Numbered newest first within each category.\par}
\nocite{*}
\pubsection{article}{J}{Journal Articles}
\pubsection{inproceedings}{C}{Conference Proceedings}
\pubsection{patent}{P}{Patents}
\pubsection{unpublished}{X}{Preprints and Manuscripts}
